% Elementos comuns aos diagramas de sequência (uml_borda.tex e uml_nuvem.tex).
% Usado pelo texto (tcc2.tex) e pela apresentação (apresentacao_tcc2.tex).
% Requer as bibliotecas TikZ arrows.meta e calc.
%
% Posição horizontal de cada linha de vida (cm)
\def\umlXno{0}
\def\umlXrad{6.0}
\def\umlXrx{10.5}
\def\umlXsrv{15.0}
%
% \umlcabecalho{x}{texto}: caixa do participante no topo
\providecommand{\umlcabecalho}[2]{%
  \node[draw, fill=white, rounded corners=1pt, align=center, minimum width=3.2cm,
        minimum height=0.9cm, font=\small] at (#1,0) {#2};}
% \umlvida{x}{y final}: linha de vida tracejada
\providecommand{\umlvida}[2]{\draw[dashed, black!55] (#1,-0.45) -- (#1,#2);}
% \umlmsg{y}{x origem}{x destino}{rótulo}: mensagem síncrona
\providecommand{\umlmsg}[4]{%
  \draw[-{Stealth[length=2.2mm]}, thick] (#2,#1) -- node[above, font=\footnotesize,
    align=center, inner sep=2pt] {#4} (#3,#1);}
% \umlret{y}{x origem}{x destino}{rótulo}: resposta (tracejada)
\providecommand{\umlret}[4]{%
  \draw[-{Stealth[length=2.2mm, open]}, dashed] (#2,#1) -- node[above, font=\footnotesize,
    align=center, inner sep=2pt] {#4} (#3,#1);}
% \umlauto{y}{x}{rótulo}: mensagem para o próprio participante
\providecommand{\umlauto}[3]{%
  \draw[-{Stealth[length=2.2mm]}, thick] (#2,#1) -- ++(0.45,0) -- ++(0,-0.3) -- ++(-0.45,0);
  \node[anchor=west, font=\footnotesize, align=left, fill=white, inner sep=1pt]
       at (#2+0.55,#1-0.15) {#3};}
% \umlloop{y topo}{y base}{x esq}{x dir}{condição}: fragmento combinado "loop"
\providecommand{\umlloop}[5]{%
  \draw[black!70] (#3,#1) rectangle (#4,#2);
  \node[draw, fill=white, anchor=north west, font=\footnotesize\bfseries, inner sep=2pt]
       (umlL) at (#3,#1) {loop};
  \node[anchor=west, font=\footnotesize] at (umlL.east) {[#5]};}
